\section{Reproduksi dan Pengujian}
\label{sec:reproduksi}

\subsection{Lingkungan eksekusi}
Program menggunakan Python 3.10 atau lebih baru. Eksekusi yang diperiksa
menggunakan Python 3.10, NumPy 2.2.6, SciPy 1.15.3, scikit-image 0.25.2,
Matplotlib 3.10.9, dan Pillow 12.3.0. Semua dependensi program dicantumkan
di \path{requirements.txt}; dependensi pengujian tambahan tercantum pada
\path{requirements-dev.txt}.

Satu cara eksekusi digunakan secara konsisten, yaitu melalui Python di
\kode{.venv}. Dari folder utama repositori pada Windows PowerShell,
persiapan berikut cukup dilakukan sekali:
\begin{lstlisting}[style=terminal]
python -m venv .venv
.\.venv\Scripts\python.exe -m pip install -r requirements.txt
\end{lstlisting}
Setelah dependensi terpasang, jalankan ketiga program:
\begin{lstlisting}[style=terminal]
.\.venv\Scripts\python.exe soal1/main.py
Push-Location soal2
..\.venv\Scripts\python.exe main.py
Pop-Location
.\.venv\Scripts\python.exe soal3/main.py
\end{lstlisting}
Pemanggilan langsung interpreter tersebut tidak memerlukan aktivasi
lingkungan. Folder \kode{.venv} dibuat pada masing-masing komputer dan
diabaikan oleh Git.

Nomor 2 dijalankan dari folder \kode{soal2/} karena lokasi masukan
dan keluaran mengikuti folder kerja. Tutup jendela visualisasi agar terminal
dapat melanjutkan. Kode pekerjaan Dustin dipertahankan.

\subsection{Berkas keluaran}
Hasil nomor 1 tersimpan di \path{soal1/hasil/}. Berkas \kode{im1.npy},
\kode{im2.npy}, dan \kode{hasil.npy} menyimpan data numerik bertanda.
Berkas PNG merupakan visualisasi, sedangkan \kode{metadata.json}
menyimpan parameter dan statistik respons.

Hasil nomor 2 tersimpan di \path{soal2/hasil_karakter/}: 18 gambar
karakter serta \kode{_visualisasi_bbox.png}. Pemeriksaan eksekusi dan hasil
nomor 2 dijelaskan pada Bagian~\ref{sec:soal2}; tidak termasuk dalam
tes otomatis nomor 1 dan 3.

Hasil nomor 3 tersimpan di \path{soal3/hasil/}. Berkas \kode{mask_koin.png}
menunjukkan segmentasi; \kode{label_koin.png} membedakan objek;
\kode{anotasi_koin.png} menunjukkan ID dan kelompok. Pengukuran tiap objek
tersimpan di \kode{objek_koin.csv}, sedangkan jumlah serta statistik setiap
kelompok tersimpan di \kode{ringkasan.json}.

\subsection{Pemeriksaan otomatis}
Dua belas pengujian mencakup tiga pemeriksaan LoG, lima pemeriksaan koin
inti, dan empat kasus pengembangan watershed. Pengujian dilakukan dengan:
\begin{lstlisting}[style=terminal]
.\.venv\Scripts\python.exe -m pip install -r requirements-dev.txt
.\.venv\Scripts\python.exe -m pytest -q
\end{lstlisting}

\begin{table}[H]
\centering
\caption{Cakupan dan hasil delapan pengujian inti.}
\label{tab:pengujian}
\small
\begin{tabularx}{\textwidth}{@{}p{0.7cm}Xp{1.2cm}@{}}
\toprule
\textbf{No.} & \textbf{Pemeriksaan} & \textbf{Hasil}\\
\midrule
1 & LoG pada citra konstan $0{,}4$ dengan $\sigma=0{,}1$ menghasilkan $-80$. & Lulus\\
2 & Dua respons dengan sigma sama mempunyai selisih nol persis. & Lulus\\
3 & Impuls pada sigma 1 dan 2 menghasilkan respons bertanda; penyimpanan
\kode{.npy} mempertahankan pengurangan. & Lulus\\
4 & Seluruh 100 objek dipasangkan ke pusat acuan; diameter dan kelompok
ukurannya sesuai toleransi pemeriksaan. & Lulus\\
5 & Inversi intensitas, sehingga latar menjadi terang, mempertahankan
deteksi dan kelompok seluruh koin. & Lulus\\
6 & Resolusi setengah mempertahankan jumlah dan kelas; diameter mengikuti
perubahan skala. & Lulus\\
7 & Lubang diisi dan derau kecil dibuang pada contoh sembilan koin; hasil
dua kelompok berjumlah tujuh dan dua. & Lulus\\
8 & Citra seragam dan berkas citra kosong ditolak dengan pesan kesalahan. & Lulus\\
\bottomrule
\end{tabularx}
\end{table}

Empat kasus tambahan mencakup kesesuaian mask pada gradien latar dengan
dua polaritas, pemetaan satu per satu terhadap 24 pusat koin foto nyata,
dan penolakan penanda watershed tidak valid atau kosong. Pada citra
gradien sintetis, rasio irisan terhadap gabungan mask acuan lebih dari
$0{,}94$. Hasil seluruh pengujian adalah \texttt{12 passed}.
Pengujian pada sigma 1 dan 2
merupakan pemeriksaan tambahan terhadap penyimpanan nilai bertanda;
hasil utama nomor 1 tetap memakai sigma $0{,}1$ dan $0{,}2$.
Contoh sembilan koin menunjukkan bahwa jumlah hasil tidak dipaksa menjadi
100. Keberhasilan pada data sintetis tidak menyatakan akurasi universal
pada foto nyata dengan pencahayaan, perspektif, atau latar yang berbeda.
